% =============================================================================
% 第7章：结论
% =============================================================================

\section{结论}
\label{sec:conclusion}

本文以FT-Transformer为载体，对自注意力机制在PCA匿名化表格数据上的失效模式进行了系统诊断。通过四个变体的对照实验，考察了两条修复路径的精确边界及其交互效应。全文凝练为三个递进的结论。

\subsection{结论一：PCA空间是注意力塌缩的根因}

嵌入SVD分析表明，29个特征专属的64维嵌入向量几乎全部位于3--4个主方向张成的低维子空间中，有效秩远低于名义维度。PCA匿名化使所有V特征的数学属性高度同构——它们的数值分布相似，"身份"仅由在主成分空间中的坐标位置区分。Feature Tokenizer从几乎相同的输入分布中学习截然不同的嵌入方向，面临梯度层面的区分性困难。这一结构性前提使自注意力在PCA匿名化数据上系统性失效——无论后续如何修复归一化函数，嵌入层都缺乏为注意力提供方向鉴别力的能力。

\subsection{结论二：两条修复路径各有边界，联合使用存在结构性拮抗}

归一化端（Sparsemax）恢复注意力多样性$\times$7.5，AUPRC提升至0.9384（+2.6pp），ECE改善至0.0120（$\sim$20倍），但修复的是Softmax饱和的"症状"而非嵌入同构的"病因"。数据端（LOR-MIM）通过门控混合打破PCA对称，实现最大多样性恢复（$\times$11.2--16.5），但AUPRC改善有限（+1.2pp），且与Sparsemax联合使用时产生拮抗效应——LOR-MIM+Sparsemax的注意力多样性达$\times$16.5但AUPRC仅为0.9267，低于Sparsemax单独使用的0.9384。LOR-MIM的混合破坏了Sparsemax稀疏选择所依赖的logits序关系：数据端修复让注意力"活过来"，但归一化端修复的"选择"因此偏离方向。

这一拮抗效应构成了本文最核心的科学发现：在PCA空间中，"注意力能不能区分特征"和"注意力区分得对不对"是两个相互独立且可能冲突的维度。两条修复路径各自提升其中一个维度，但无法同时兼顾两者——这正是未来工作的明确方向。

\subsection{结论三：诊断优于比较}

本文隐含的方法论立场是：系统\textbf{诊断}失效模式比\textbf{比较}性能数字提供更多可迁移的知识。四个变体的AUPRC数值无法揭示拮抗关系——是LOR-MIM+Sparsemax联合实验的设计，将两条修复路径从"独立有效"重新框定为"相互拮抗"。如果论文的核心叙事是"方法A优于方法B"，其结论受限于数据集；如果是"方法A与B在机制X上拮抗"，其结论是可迁移的。

\subsection{对金融行业的实践启示}

本文的诊断发现对金融机构的模型选型具有直接的操作指导意义：

\begin{enumerate}
    \item \textbf{在使用PCA匿名化数据时，优先验证注意力塌缩风险}。建议金融机构在使用Transformer类模型处理匿名化表格数据前，先执行本文提出的注意力分布检验（计算$\sigma/(1/d)$比值）——若该值低于0.05，塌缩风险较高，需考虑替代方案。
    \item \textbf{根据数据匿名化程度选择模型架构}。数据经PCA匿名化且无任何语义锚点时，Sparsemax替代Softmax是一个低成本的改善方案（AUPRC +2.6pp，改一行代码）。若业务允许接受BCE在不平衡场景下的理论风险，在本实验条件下BCE训练提供了最高的排序精度（AUPRC 0.9491）。LOR-MIM适合作为"附加诊断层"——将其加入pipeline后注意力多样性的大幅提升可作为模型是否对特征结构敏感的检验信号。
    \item \textbf{将注意力多样性指标纳入风控模型上线前的质量检查清单}。塌缩的注意力意味着模型无法解释"为什么这笔交易被判定为欺诈"——在面临监管审查时，均匀的注意力分布无法提供任何特征归因信息。定期监测注意力分布，可在模型性能尚未出现显著下降之前发现特征匿名化带来的结构性退化。
\end{enumerate}

\textbf{局限与未来工作。}本文的局限——单数据集、单种子（seed=42）、诊断框架的外部验证缺失——已在正文中指出。PCA匿名化场景比表面看起来更普遍：联邦学习中的差分隐私梯度扰动、多传感器阵列的同质测量、差分隐私表格数据发布中的随机投影，都可能产生类似的"语义真空"效应。本文的发现和诊断框架在这些场景中的适用性，构成未来研究的重要方向。
